\section{WWW 2026 研究主题分析}

\writingguide{建议写 1200--1500 字。先说明论文样本来自哪里、如何筛选，再对论文进行主题分类。每类至少选取若干论文作比较，避免“一篇论文一段话”的流水账。}

\subsection{论文样本与分类方法}

本文以 WWW 2026 论文集及 Track 页面列出的论文为样本，根据研究问题、输入模态、检索对象和评价目标进行归类。

\placeholder{填写样本数量、筛选日期、纳入与排除标准；核对每篇论文确实属于本 Track。}

\begin{table}[htbp]
  \centering
  \caption{Track 论文的主题分类示例}
  \label{tab:topic-taxonomy}
  \begin{tabularx}{\textwidth}{>{\raggedright\arraybackslash}p{0.18\textwidth} X X}
    \toprule
    主题类别 & 主要研究问题 & 代表论文 \\
    \midrule
    检索与重排序 & 如何提高复杂查询下的召回率、相关性与效率 & \placeholder{论文 A、论文 B} \\
    检索增强生成 & 如何选择可靠证据并降低生成幻觉 & \placeholder{论文 C、论文 D} \\
    多模态检索 & 如何对齐文本、图像、视频等异构信息 & \placeholder{论文 E、论文 F} \\
    智能体与工具检索 & 如何为 Web Agent 选择工具、记忆与行动依据 & \placeholder{论文 G、论文 H} \\
    \bottomrule
  \end{tabularx}
\end{table}

\subsection{主要研究主题}

\placeholder{按照表~\ref{tab:topic-taxonomy} 的分类逐类分析共同问题、代表方法和实验结论。}

\subsection{横向比较与阶段性特征}

\placeholder{比较不同研究路线对准确性、实时性、可解释性、计算成本与安全性的权衡，并总结 WWW 2026 反映出的整体趋势。}
